\documentclass[11pt]{article}

\usepackage[utf8]{inputenc}
\usepackage[T1]{fontenc}
\usepackage{mathptmx}
\usepackage[margin=2.3cm]{geometry}
\usepackage{microtype}
\usepackage{enumitem}
\usepackage{titlesec}
\usepackage{parskip}

\titleformat{\section}{\normalfont\bfseries\large}{}{0pt}{}
\titlespacing*{\section}{0pt}{10pt}{4pt}
\setlist[itemize]{leftmargin=1.2em,itemsep=2pt,topsep=2pt}
\pagestyle{empty}
\setlength{\parskip}{5pt}
\hyphenpenalty=10000
\exhyphenpenalty=10000
\emergencystretch=2em

\begin{document}

\begin{center}
	{\small SMAII, Redes Complejas Multicapa}\\[6pt]
	{\Large\bfseries Propuesta de proyecto semestral}\\[3pt]
	{\normalsize Efecto de una campaña sobre la opinión colectiva y predicción de su conveniencia}\\[4pt]
	{\small Septiembre de 2026}
\end{center}

\section*{Integrantes del equipo}
\begin{itemize}
    \item Jorge Miyoshi Mikami Alonso
    \item Luis Manuel Reyes Colín 
\end{itemize}

\section*{Planteamiento}

El proyecto parte de una decisión concreta que enfrenta cualquier emisor de un mensaje
masivo, ya sea una campaña de marketing, un anuncio político o una publicación en redes.
La intención suele ser desplazar la opinión del público hacia una posición determinada,
pero el resultado no siempre acompaña a la intención. El mensaje puede lograr su objetivo,
puede resultar irrelevante o puede provocar el efecto contrario y consolidar al grupo que
se le opone. Nuestro interés es determinar si ese desenlace puede anticiparse antes de
que la campaña se lance.

Para abordarlo simularemos una población conectada en la que cada individuo sostiene una
opinión y solo se deja influir por quienes piensan de manera suficientemente parecida.
Sobre esa base incorporaremos el estímulo externo que representa la campaña y el sesgo
que introducen las plataformas digitales al favorecer el contacto entre personas
afines. Analizaremos hacia dónde evoluciona el sistema según la apertura de la población,
la intensidad del sesgo, la frecuencia con la que se repite el mensaje y qué tan extrema
es la posición que promueve.

El resultado que buscamos es una caracterización de las condiciones bajo las cuales
conviene lanzar la campaña y aquellas en las que resulta desaconsejable. Consideramos
también una segunda capa en la que la opinión se traduzca en una conducta observable,
como una compra o un voto, dado que ese es el desenlace que finalmente interesa al emisor.

\section*{Metodología}

El trabajo se apoyará principalmente en simulación sobre redes, primero sintéticas y
después, en la medida en que consigamos datos adecuados, sobre una red real. Recorreremos
distintas combinaciones de parámetros y evaluaremos el estado final de la población,
distinguiendo entre consenso, polarización y fragmentación.

Dado el volumen de escenarios que esto genera, planeamos apoyarnos en aprendizaje
automático para clasificarlos y para entrenar un modelo que, a partir de las condiciones
iniciales y las características del mensaje, prediga si la campaña resultará favorable o
contraproducente. Este sería el principal entregable del proyecto.

\section*{Referencias de partida}

\begin{itemize}
	\item V. Pansanella, A. Sîrbu, J. Kertesz and G. Rossetti, 2023.
	      \emph{Mass media impact on opinion evolution in biased digital environments: a bounded confidence model}.
	      Scientific Reports \textbf{13}, 14600.

	\item L. Yin, J. Chen, F. Gao and X. Wu, 2025.
	      \emph{Coevolution of opinion and consumption behavior under a two-layer network framework}.
	      Physica A \textbf{680}, 131020.
\end{itemize}
\end{document}
